\section{The dyadic count and the explicit Cauchy modulus}
\label{sec:convergence}

\subsection{A uniform numerical bound}

\begin{lemma}
\label{lem:polynomial}
For every \(t\ge0\), \((t+1)^2\le2^{2t}\).
\end{lemma}

\begin{proof}
At \(t=0\) both sides are \(1\).
Since \(t+2\le2(t+1)\), squaring and using the induction hypothesis gives
\[
 (t+2)^2\le4(t+1)^2\le4\cdot2^{2t}=2^{2(t+1)}.
\]
This is induction on a quantifier-free numerical inequality.
\end{proof}

\begin{proposition}[Explicit dyadic counting theorem]
\label{prop:dyadic-count}
For every \(k\ge0\),
\[
 k^2T(2^k)\le1400\,2^k\le2048\,2^k.
\]
\end{proposition}

\begin{proof}
Put
\[
 X=2^k,\qquad t=k\mathbin{\mathrm{div}}10,\qquad z=2^{2t}.
\]
Division with remainder gives
\[
 10t\le k<10(t+1),\qquad
 k^2\le100(t+1)^2.
\]
Lemma~\ref{lem:harmonic-dyadic} shows that the main term in
Proposition~\ref{prop:finite-sieve} satisfies
\[
 \frac{k^2X}{G(z)}
 \le\frac{12k^2X}{(t+2)^2}
 \le1200\,\frac{(t+1)^2}{(t+2)^2}\,X
 \le1200X.
\]
For the error, Lemma~\ref{lem:polynomial} gives
\[
 k^2z^4\le100(t+1)^2\,2^{8t}
          \le100\,2^{10t}
          \le100\,2^k=100X.
\]
Since \(z\ge1\), also \(z\le z^4\), and so \(k^2z\le100X\).
Multiply \eqref{eq:finite-sieve} by \(k^2\) and add the three bounds:
\[
 k^2T(2^k)
 \le k^2z+\frac{k^2X}{G(z)}+k^2z^4
 \le(100+1200+100)X=1400X.
\]
Finally \(1400\le2048=2^{11}\).
This proof includes \(t=0\), and even \(k=0\); it never divides by \(k\).
For \(k\ge1\), division by \(k^2\) yields the familiar form
\[
 T(2^k)\le2048\,\frac{2^k}{k^2}.
\]
\end{proof}

\subsection{A finite tail estimate}

Write \(C=2048\). For a dyadic shell indexed by \(j\ge1\), put
\[
 V_j=\sum_{\substack{2^j\le p<2^{j+1}\\\Twin(p)}}
              \left(\frac1p+\frac1{p+2}\right).
\]
Each summand is at most \(2/p\le2/2^j\).
The number of summands is at most \(T(2^{j+1})\). The inclusive endpoint
in \(T\) only enlarges this bound. Proposition~\ref{prop:dyadic-count}
therefore gives
\begin{equation}
 V_j\le\frac2{2^j}T(2^{j+1})
       \le\frac{4C}{(j+1)^2}.
 \label{eq:shell}
\end{equation}

For \(j\ge1\), all denominators below are positive, and
\begin{equation}
 \frac1{(j+1)^2}\le\frac1{j(j+1)}
  =\frac1j-\frac1{j+1}.
 \label{eq:telescope-step}
\end{equation}
Summing only from \(j=m\) to a finite \(K\ge m\) gives
\begin{equation}
 \sum_{j=m}^{K}\frac1{(j+1)^2}
 \le\frac1m-\frac1{K+1}\le\frac1m.
 \label{eq:telescope}
\end{equation}
This uses no assertion about the convergence of an auxiliary infinite
series.

\begin{proposition}[Finite reciprocal tail]
\label{prop:tail}
For \(m\ge1\) and \(L\ge0\),
\[
 \sum_{\substack{2^m\le p<L\\\Twin(p)}}
          \left(\frac1p+\frac1{p+2}\right)
 \le\frac{4C}{m}=\frac{8192}{m}.
\]
\end{proposition}

\begin{proof}
If \(L\le2^m\), the sum is empty. Otherwise choose, for instance,
\(K=m+L\). The elementary inequality \(L<2^{K+1}\) follows from
\(u<2^{u+1}\), proved by quantifier-free induction.
Every integer \(p\) with \(2^m\le p<L\) lies in a unique shell
\([2^j,2^{j+1})\) with \(m\le j\le K\).
To find its index, use the bounded search for the first
\(j\in\{m,\ldots,K\}\) with \(p<2^{j+1}\); the previous endpoint,
or \(p\ge2^m\) in the first case, gives \(2^j\le p\).
Thus the original sum is at most \(\sum_{j=m}^K V_j\).
Equations~\eqref{eq:shell} and \eqref{eq:telescope} give the claimed bound.
\end{proof}

\subsection{Convergence with a primitive recursive modulus}

Fix \(q\ge0\), and set
\[
 m=2^{q+13},\qquad N(q)=2^m.
\]
Then \(m\ge1\) and
\begin{equation}
 \frac{8192}{m}
 =\frac{2^{13}}{2^{q+13}}
 =2^{-q}.
 \label{eq:modulus-arithmetic}
\end{equation}
Let \(M,L\ge N(q)\). By decidable comparison, assume first that \(L\le M\).
Nonnegativity of the summands and Proposition~\ref{prop:tail} give
\[
 0\le B(M)-B(L)
 =\sum_{\substack{L\le p<M\\\Twin(p)}}
       \left(\frac1p+\frac1{p+2}\right)
 \le\sum_{\substack{2^m\le p<M\\\Twin(p)}}
       \left(\frac1p+\frac1{p+2}\right)
 \le2^{-q}.
\]
If \(M<L\), apply the same calculation with the names interchanged.
This proves Theorem~\ref{thm:main}(iii).

For the precise integer version, take \(N(q)\le U\le V\), use
\eqref{eq:tail-code}, and multiply the tail inequality by
\(2^qD(V)>0\). The result is exactly \eqref{eq:ipra-target}.
This completes the finite mathematical argument.

\begin{remark}[Effectivity and numerical usefulness]
The formula \(N(q)=2^{2^{q+13}}\) is a composition of primitive recursive
exponentiation functions. Composing it with bounded primality testing,
factorial, and finite summation remains primitive recursive.
Its size is immaterial to primitive recursiveness, although it makes this
particular approximation algorithm impractical.
The theorem supplies a certified convergence rate, not a competitive
method for computing many digits of Brun's constant.
\end{remark}
